\section{总结与延伸}

\subsection{全讲核心结论}
\begin{table}[H]
\centering
\begin{tabular}{p{3.0cm}p{4.0cm}p{4.4cm}}
\toprule
\textbf{主题} & \textbf{关键结论} & \textbf{对实践的直接影响} \\
\midrule
系统演进主线 & AI 工作负载迫使基础设施重写 & 需要拓扑感知调度与数据面重构 \\
浪潮切换 & RAG 等能力从显式热点走向隐式基础设施 & 重点从概念竞争转向系统集成能力 \\
生产可靠性 & 故障恢复阶段最容易触发级联崩溃 & 必须建设流控、预热、降级、演练机制 \\
硬件软件关系 & 高密度节点与高带宽互联成为主流 & 平台抽象需暴露关键物理约束 \\
创业机会 & 垂直企业场景仍有大量未解问题 & 结合行业知识做高价值差异化交付 \\
未来方向 & 中心化大模型与边缘高效推理并存 & 双路线并行：上限探索 + 效率利用 \\
\bottomrule
\end{tabular}
\caption{Lecture 10 的系统设计结论与工程含义}
\end{table}

\begin{importantbox}{一句话总结}
Agent 时代的竞争，不再只是 ``谁的模型更强''，而是 ``谁能把模型能力稳定、低成本、可治理地变成真实系统能力''。
\end{importantbox}

\subsection{进一步阅读}
\begin{itemize}
    \item Berkeley RDI Agentic AI MOOC F25 其余讲次（安全、评测、多智能体、生产部署）。
    \item Snowflake 与 Databricks 的 Data Cloud 架构公开资料。
    \item NVIDIA DGX 平台与大规模训练系统工程文档。
    \item 云原生与 AI 基础设施融合实践（Kubernetes, Slurm, Ray, SkyPilot 等）。
    \item 企业级 RAG/Agent 评测方法（事实性、可追溯性、成本与时延联合评估）。
\end{itemize}

\end{document}
